% !TeX root = ../papers/nc-p3.tex
\section{各矢に高々二重みのある全ての乗法群作用}
\label{sec:binary-torus-windows}
\begin{theorem}\label{thm:binary-torus-windows}
標準正則窓の非自明な乗法群が各矢に高々二つの重みを持つとする。各位置のスカラーを除くと
四位置とも同じ正の重み差を持ち、重み重複度も同じである。共通の重み差を1にする。
\begin{enumerate}
\item 重複度 $3+1$（または反転）の全窓は、三生成二次三次元Koszul左右AS正則な通常の
次数付き代数 $S$ に中心文字を加えた $S[z]$ の添字反復であり、最小周期は1である。
重みを保つ添字同型は $S$ の次数付き同型で完全に決まる。
\item 重複度 $2+2$ の全窓は定理\ref{thm:split-inverse-transporter}の可逆チャートで尽きる。
特に二軸の純二次関係のどれか一つでも退化すれば正則でない。
定理\ref{thm:split-jordan-transport}と系\ref{cor:split-finite-unmarking}の
スペクトル条件・残余共役・有限個の印の検査はこの表現型全体に適用できる。
全て二次PBWで、多重安定極限も同じチャートで正則である。
\end{enumerate}
\end{theorem}

\begin{proof}[等間隔の不均一型の排除]
まず全ての非自明な位置で差が1の場合を扱う。系\ref{cor:two-weight-profiles}の $0121$ では
位置1の高位文字の双対点で縮約する。位置0の端支持は $B_0$ の重み1の次元から2以下であり、
位置2の端支持は最上重み1なので二次元の低位空間に入る。
定理\ref{thm:all-endpoint-ranks-two}によって両端が2となり、定理\ref{thm:no-two-sided-centers}に反する。
$0323$ では重みを反転しスカラー位置を正規化すれば同じである。

$1133$ の $B_1=(0,6,0)$ の場合、全ての台で位置1,2の高位文字の個数は1である。
従ってこの二位置だけで高位を拡大する別の乗法群が窓を保つ。それは隣接二位置だけに非自明で、
すでに系\ref{cor:two-weight-profiles}のどの行にもない。
残る $B_0=(4,2,0),B_1=(1,4,1)$ では、位置0,1の高位直線を $kh_0,kh_1$、
位置2,3の低位直線を $kx_2,kx_3$ とし、その補三空間を $W_i$ とする。
$01|23$ の混合切断は二つの三空間ずつの間の階数2行列で、各側の像は二平面である。
$12|30$ の重み0と2が各1なので、$R_1$ に $a_1x_2,h_1a_2$、
$R_3$ に $a_3h_0,x_3a_0$ という純関係がある（$0\ne a_i\in W_i$）。
従って $R_0$ の混合二平面は $h_0W_1,W_0h_1$ の一方だけに属する非零元を持てない。
どちらも上の純関係と連続するためである。同じく $R_2$ の二平面にも $x_2 W_3,W_2x_3$ の単独支持はない。
よって四つの三空間への射影は全て階数2で、混合切断の各三次正方ブロックは階数2となる。
しかし $h_0W_1x_2W_3$ 成分は $12|30$ の重み0のブロックそのもので階数1である。これも矛盾する。
\end{proof}

\begin{proof}[異なる重み差の排除]
非自明な位置数 $n\le4$ とし、高位・低位の選択を立方体 $\{0,1\}^n$ の頂点で表す。
実際に係数のある頂点は、正の差の列 $d$ と最上重み $H$ に対する $d\cdot x=H$ に含まれる。
この台上で定数となる任意の整数係数の一次式も射影窓を保つ一径数を与える。
従って台のアフィン包に直交する非零の $\{-1,0,1\}$ ベクトルで、
四位置へ拡張したとき零成分を持つものはない。あれば負符号の位置では二重みを取り替えてスカラーを引き、
前段で排除したスカラー位置付きの等間隔作用になるからである。

台のアフィン基底として高々 $n$ 頂点を取れる。その差ベクトルの直交を全て調べると、上記条件に適い
正の法線を持つ基底は $n=4$ の四点のみで、原始法線の成分は等しいか、並べ替えて $(1,1,1,2)$ である。
後者の $H$ は2または3で、反転により2とできる。巡回で差2を最後に置き、
$k_i\in\{1,2,3\}$ に対して $\sum d_i k_i=8$ を課す。各隣接対と相補対の重み次元から
六関係の重複度を列挙する。定理\ref{thm:torus-euler-characters}の $p_1$ の式に出る重みは0から5に入り、
四指数の和は8である。全候補が $(t,y)=(11,7)$ でその四一次因子の積と不一致となる。
二つの有限整数計算は \path{computations/verify_binary_torus_windows.py} による。
\end{proof}

\begin{proof}[純粋関係と曲面、動く軌道の終点]
以後 $X_i,Y_i$ を低位・高位空間とする。曲面射影を持つ正則窓には純粋関係 $u v$ はない。
定理\ref{thm:no-moving-plane-pencils}と定理\ref{thm:isolated-surface-fibers}によれば、
整二次曲面商では全空間を張る整なグラフの一般点でこの積は消えない。
平面商では正則中心一次拡大の表示で、中心文字による随伴次数は三変数多項式環の中心多項式拡大であり整域である。
多項式環自身でも明らかである。従って以下で純関係が生じれば、全射影像が高々曲線である。

曲面のない場合に、乗法群が非自明に作用する後続直線を持つ固定点から
定理\ref{thm:finite-branch-corridors}の曲線列をたどる。
一般後続は一意なので全ての中間写像は同変である。
各中間像は非自明軌道の閉包で、二重みしかないため高位軸の一点と低位軸の一点を結ぶ直線である。
最後の直線では全点の後続がただ一つの同じ終点になる
（逆向きに同定理を適用、または命題\ref{prop:branch-corridor-degree}の $f_m=0$）。
特にこの直線の二つの軸端点にも同じことがいえる。
\end{proof}

\begin{proof}[$3+1$ の正規直線]
高位文字を $z_i$ とし、その双対点を $p_i$ とする。$B_i=(3,3,0)$ なので全 $p_i$ は連続して点列となる。
混合三関係は $X_i z_{i+1}\oplus z_i X_{i+1}$ の三平面である。
一方の射影が同型でないと純関係があり、対応する $p_j$ での前または後の評価階数が2以下である。
時間を必要なら逆にして $p_j$ の後続が直線となる場合を扱う。
その直線は $p_{j+1}$ を含み、高位空間は一次元なので群が非自明に作用する。
前段の曲線列の各中間直線の高位端はその位置の $p_i$ である。
最後の高位端には元の全高位点列で高位の後続があるから、列の定終点も高位点である。
従って固定した二つの高位点の間に、最初の直線の点を変えた異なる経路がある。
高位点列で同じ開始位置へ四周期を合わせて戻れば、二経路を独立に反復して指数増大する。
実際異なる内部点での評価は非比例で、両端を固定した各ブロックの独立な評価をテンソル積にできる。
これは標準成長に反する。

よって全ての混合射影は同型で、$X_i z_{i+1}=z_i X_{i+1}$ となる。
$z$ を右へ動かすと、全語は低位語の後に $z$ の列を付けたもので張られる。
重み0の部分代数の各長さの次元はKoszul指標の逆公式で $\binom{n+2}{2}$ である。
これは各位置の重み0のKoszul係数が $(1,3,3,1)$ だからである。
全長 $n$ のこれらの上界を足すとちょうど $\binom{n+3}{3}$ になるため各正規表示は独立である。
逆向きにも同じで $z$ 列は正則正規となる。定理\ref{thm:regular-normal-linear}で $z$ を中心化する
位置移動は低位空間を保つ。低位関係は純粋なので補正 $\alpha,\beta$ は零であり、$S[z]$ を得る。
逆も同定理による。印を保つ同型は中心化と可換で、元の $S$ の次数付き同型を与える。
三生成関係の全次数判定には補題\ref{lem:ternary-finite-graph-test}を使える。
\end{proof}

\begin{proof}[$2+2$ の特異純次数関係の排除]
色階数は $(1,4,1)$ で、各軸に一本ずつの関係がある。どの軸上の点も同じ軸上に少なくとも一つの後続を持つ。
従って非自明に動く後続直線から始まる上記の列は存在できない。
最後の直線の二軸端点が、それぞれ自分と同じ軸に後続を持ち、定終点が互いに素な二つの軸に同時に属するからである。

低位の関係 $r_i$ が特異と仮定する。これは純関係で、曲面はない。$r_i$ を左で消す低位双対点 $p$ は、
低位の順軸直線全体を後続に持つ。一段後のその直線から始めて一般後続を取る際、
次の低位関係 $r_{i+1}$ も特異なら一般点の低位後続は一定である。
一般後続は一意なので動く位置が一個だけとなり、上記の列の $m\ge2$ に反する。
従って $r_{i+1}$ は可逆であり、次の動く像も低位の直線で $d_1=d_2=1$ である。
命題\ref{prop:branch-corridor-degree}の次数式の最初の項から $f_1=3-d_2=2$。
ゆえに最初の低位直線上に順側評価が階数2の点がある。そこでの純低位後続は $r_{i+1}$ の可逆性から一点のみである。
従ってその後続直線は不変で低位・高位をまたぎ、群が非自明に作用する。
これは前段で排除した場合である。高位の関係も重みの反転で同様であり、
定理\ref{thm:split-inverse-regularity}と定理\ref{thm:split-polystable-good-limits}が結論を与える。
\end{proof}

\begin{corollary}\label{cor:binary-torus-unmark}
二つの表現型の内部と相互に、印を外した同型は有限個の重み付き記録の集合の一致で決定される。
\end{corollary}
\begin{proof}
窓の射影安定群の一つの極大トーラス内で、上のいずれかの正規化された射影重みを持つ全ての一径数を探す。
各矢の固有文字の任意の二つの差への値は $0,1,-1$ に限られ、これらの差は極大トーラスの
有理指標空間を張る（射影四位置への作用が忠実）。従って候補は有限である。
全ての印は共役によりここに入り、共役な印は同じ重み付き記録を与える。
逆に一つでも記録の一致があれば、重みを保つ同型が実際の窓の同型である。
位置を巡回した窓との比較にも同じ判定を用いる。
\end{proof}
